\textbf{Selección del motor}

Una vez calculada la potencia consumida, por el eje de aireación es posible calcular y seleccionar el motor necesario para mover la carga.

Dado que se requiere una velocidad comprendida entre 3 y 6 rpm, y una potencia baja, el motor a seleccionar es un motor de baja potencia y baja velocidad.
$monofásico$

\begin{table}[h]
    \centering
    \caption{Características del motor seleccionado} 
    \begin{tabular}{c|c|c}
        Potencia & Tensión & RPM \\
        1/4 hp  & 127-220 V & 1750 
    \end{tabular}
    \label{tab:DatosMotro}
\end{table}

Se propone un reductor de velocidad conectado a la flecha del motor principal, el reductor se compone de una corona y un tornillo sin fin, ya que este mecanismo permite una alta relación de reducción.

La segunda etapa de reducción está compuesta por una banda de transmisión, ya que al ser un un mecanismo síncrono permite una velocidad constante a la salida, además de una eficiencia del $95\%$.

Al considerar las perdidas de potencia en las distintas etapas de reducción, se calculó la potencia real requerida. 

La eficiencia del reductor de corona y tornillo sin fin es de $82\%$.

Se tomará una eficiencia de $75\%$ para la caja reductora y $90\%$ para la banda.\cite{Mott}

\begin{center}
    $P_{Req}=\frac{(25.2 W)}{(0.75)(0.9)}$\\
\end{center}
\begin{center}    
    $P_{Req}=37.33 W\%$\\
\end{center}
\begin{center}
    $P_{Req}\approx 37.5 W\%$ 
\end{center}

